% !TEX root = ../main.tex
\documentclass[../main.tex]{subfiles}

\begin{document}

\selectlanguage{english}

\begin{center}
    \textbf{\LARGE COMPREHENSIVE 5-CHAPTER INTRO TO AI EXAM SOURCE} \\
    \vspace{0.2cm}
    \textbf{\LARGE Complete Academic Framework, Mathematical Foundations \& Algorithmic Features} \\
    \textit{(Strictly Mapped with PTIT Course Syllabus Rules and Russell \& Norvig Reference Core)}
\end{center}

\vspace{0.5cm}

\section{CHAPTER 1: INTRODUCTION}

\subsection{1.1. Artificial intelligence definition and concepts}
The textbook by Russell \& Norvig categorizes historical AI definitions along a two-dimensional matrix: processes/reasoning vs. behavior, and human-like performance vs. ideal rationality.
\begin{itemize}
    \item \textit{Thinking Humanly:} Cognitive modeling frameworks that formalize the internal mechanics of the human mind via empirical psychological testing.
    \item \textit{Acting Humanly:} Passing the operational Turing Test through distinct computational capabilities: NLP, knowledge representation, automated reasoning, and machine learning.
    \item \textit{Thinking Rationally:} Codifying the ``laws of thought'' using formal syllogisms and logical notations to build irrefutable argument structures.
    \item \textit{Acting Rationally (Textbook Core):} Designing artifacts that act as rational agents, selecting the absolute highest-utility action paths to maximize performance measures under environmental constraints.
\end{itemize}

\subsection{1.2. History of AI}
The historical evolution follows distinct operational eras: Gestation (1943--1955) initiating with the binary neural model of McCulloch \& Pitts; The formal birth of AI at the Dartmouth conference (1956); Great expectations and initial breakthroughs (1952--1969); The first AI winter triggered by scale limits (1966--1973); The rise of knowledge-based Expert Systems (1969--1979); The second commercial AI winter (1984--1987); and the Modern Era (1995--Present) driven by big data, rigorous probability foundations, and advanced deep learning infrastructures.

\subsection{1.3. Main research and application areas}
The essential pillars include: Computer Vision, Natural Language Processing (NLP), Autonomous Systems and Robotics, Automated Strategic Planning, Knowledge Engineering, and Machine Learning.

\subsection{1.4. Achievements and unsolved problems}
\begin{itemize}
    \item \textit{Achievements:} Superhuman performance in complex strategy board games (Deep Blue chess solver, AlphaGo), high-level autonomous vehicle routing, precision medical diagnostics, and Large Language Models (LLMs).
    \item \textit{Unsolved Problems:} Realizing Artificial General Intelligence (AGI), resolving black-box uninterpretability via Explainable AI (XAI), mitigating informational hallucination, and overcoming exponential state-space scaling barriers inherent in NP-hard optimization tasks.
\end{itemize}

\end{document}
